\documentclass[a4paper,12pt]{article}
\usepackage[T2A]{fontenc}
\usepackage[utf8]{inputenc}
\usepackage[russian]{babel}
\usepackage{amsmath,amssymb,mathtools}
\usepackage{helvet}
\renewcommand{\sfdefault}{cmss}
\renewcommand{\familydefault}{\sfdefault}
\usepackage{icomma}
\usepackage{geometry}
\usepackage{microtype}
\usepackage{tikz}
\usetikzlibrary{calc,arrows.meta,positioning,shapes.geometric}
\usepackage{tabularx,booktabs}
\usepackage{enumitem}
\usepackage{parskip}
\usepackage{fancyhdr}
\usepackage{setspace}
\usepackage{xcolor}
\usepackage[hidelinks]{hyperref}

\geometry{left=19mm,right=19mm,top=21mm,bottom=20mm}
\setlength{\headheight}{25pt}
\onehalfspacing
\setlength{\parindent}{0pt}
\setlist[enumerate]{leftmargin=8mm,itemsep=7pt,topsep=4pt}
\setlist[itemize]{leftmargin=7mm,itemsep=2pt,topsep=3pt}
\definecolor{accent}{HTML}{2A6F6B}
\definecolor{darkaccent}{HTML}{174B48}
\definecolor{textgray}{HTML}{2F3437}
\definecolor{softgray}{HTML}{DCE8E6}
\color{textgray}

\pagestyle{fancy}
\fancyhf{}
\fancyhead[L]{\small Кредиты: произвольные схемы погашения}
\fancyhead[R]{\small 13.10.26}
\fancyfoot[C]{\small \thepage}
\renewcommand{\headrulewidth}{0.4pt}
\renewcommand{\headrule}{\hbox to\headwidth{\color{softgray}\leaders\hrule height \headrulewidth\hfill}}
\newcommand{\key}[1]{\textcolor{darkaccent}{\textbf{#1}}}
\newcommand{\fboxeq}[1]{\[\boxed{#1}\]}
\newcommand{\mln}{\,\text{млн руб.}}
\newcommand{\rub}{\,\text{руб.}}

\begin{document}
\thispagestyle{empty}
\begin{center}
\vspace*{30mm}
{\Large\textcolor{accent}{\textbf{Чек-лист}}}\\[10mm]
{\Huge\textbf{Кредитные задачи}}\\[7mm]
{\Large\textbf{Произвольные схемы погашения}}\\[5mm]
{\large Долг, проценты, платежи и ограничения}\\[20mm]
{\large 13.10.26}\\[29mm]
{\large Лёвин Артём Александрович}\\[4mm]
{\large\textcolor{darkaccent}{эксклюзивно для Екатерины Гнедковой}}
\end{center}
\vfill
\begin{center}
\begin{tikzpicture}[x=1cm,y=1cm]
\draw[accent!60,line width=1.1pt]
(-7,0) .. controls (-4.1,0.95) and (-1.3,-0.75) .. (1.45,0.18)
.. controls (3.65,0.85) and (5.25,-0.5) .. (7,0.25);
\end{tikzpicture}
\end{center}
\clearpage

\section*{1. Из чего складывается каждый платёж}

В задачах на \key{произвольную схему погашения} остаток долга или размер платежа задаётся особым правилом. Удобнее всего работать с таблицей, одинаковой для всех периодов.

Пусть $S$ --- первоначальная сумма кредита, $p$ --- ставка в процентах \emph{за один период начисления}, $r=\frac{p}{100}$ --- ставка в долях, $k=1+r$ --- коэффициент роста долга.

\begin{center}
\renewcommand{\arraystretch}{1.3}
\begin{tabularx}{\linewidth}{@{}>{\raggedright\arraybackslash}X >{\raggedright\arraybackslash}X@{}}
\toprule
Столбец таблицы & Как найти \\
\midrule
Долг в начале периода & $B_{i-1}$ \\
Начисленные проценты & $I_i=rB_{i-1}$ \\
Общая задолженность & $D_i=B_{i-1}+I_i=kB_{i-1}$ \\
Платёж & $P_i=D_i-B_i$ \\
Остаток после платежа & $B_i=D_i-P_i$ \\
\bottomrule
\end{tabularx}
\end{center}

\key{Основное равенство} для $i$-го периода:
\fboxeq{B_i=kB_{i-1}-P_i,\qquad B_0=S.}
Долг после платежа становится начальным долгом следующего периода. Если кредит погашен полностью, $B_n=0$.

Если по условию известны остатки как доли начальной суммы,
\[
B_{i-1}=aS,\qquad B_i=bS,
\]
то \key{платёж находится сразу}:
\fboxeq{P_i=(ka-b)S.}

\textbf{Проверяйте порядок событий.} В рассматриваемых задачах сначала начисляют проценты, затем происходит платёж. Месячную ставку применяют к месяцам, годовую --- к годам.

\section*{2. Переплата: два равноценных способа}

Пусть $T=P_1+\dots+P_n$ --- сумма всех платежей. Если кредит погашен полностью и дополнительных комиссий нет, переплата равна сумме начисленных процентов:
\[
T-S=\sum_{i=1}^{n}I_i,\qquad
q=\frac{T-S}{S}\cdot100\%=\frac{\sum_{i=1}^{n}I_i}{S}\cdot100\%.
\]
Сравнивать переплату в процентах нужно \key{с суммой выданного кредита} $S$.

\clearpage
\section*{3. Пример: график остатка задан заранее}

Кредит $S$ выдан 15 января на полгода. В конце каждого месяца долг до внесения платежа увеличивается на $5\%$. Остатки на 15 февраля, марта, апреля, мая, июня и июля равны соответственно
\[
0{,}9S;\quad 0{,}8S;\quad 0{,}7S;\quad
0{,}6S;\quad 0{,}5S;\quad 0.
\]
Найдём платежи и процентную переплату.

\begin{center}
\small
\renewcommand{\arraystretch}{1.42}
\begin{tabular}{@{}c c c c c c@{}}
\toprule
Период & Начало & Проценты & Всего & Платёж & Остаток \\
\midrule
1 & $S$ & $0{,}05S$ & $1{,}05S$ & $0{,}15S$ & $0{,}9S$ \\
2 & $0{,}9S$ & $0{,}045S$ & $0{,}945S$ & $0{,}145S$ & $0{,}8S$ \\
3 & $0{,}8S$ & $0{,}04S$ & $0{,}84S$ & $0{,}14S$ & $0{,}7S$ \\
4 & $0{,}7S$ & $0{,}035S$ & $0{,}735S$ & $0{,}135S$ & $0{,}6S$ \\
5 & $0{,}6S$ & $0{,}03S$ & $0{,}63S$ & $0{,}13S$ & $0{,}5S$ \\
6 & $0{,}5S$ & $0{,}025S$ & $0{,}525S$ & $0{,}525S$ & $0$ \\
\bottomrule
\end{tabular}
\end{center}

Например, первый платёж:
\[
P_1=1{,}05S-0{,}9S=0{,}15S.
\]
В последнем месяце остаток после платежа равен нулю, поэтому платёж составляет всю начисленную задолженность $0{,}525S$.

\key{Быстрый путь} --- сложить проценты:
\[
\sum I_i=(0{,}05+0{,}045+0{,}04+0{,}035+0{,}03+0{,}025)S
=0{,}225S.
\]
\fboxeq{q=\frac{0{,}225S}{S}\cdot100\%=22{,}5\%.}

\textbf{Самопроверка:} сумма шести платежей равна $1{,}225S$. Это на $0{,}225S$ больше суммы кредита.

\section*{4. Как переводить ограничения в математику}

\begin{itemize}
\item \key{«Каждый платёж меньше 5 млн»} означает \emph{систему} $P_i<5$ для каждого $i$ (если платежи измерены в миллионах рублей).
\item \key{«Все платежи вместе меньше 15 млн»} означает \emph{одно} неравенство $P_1+\dots+P_n<15$.
\item \key{«Остаток на 70\% меньше прежнего»} означает $B_{i}=0{,}30B_{i-1}$, то есть сохранилось $30\%$ прежней суммы.
\end{itemize}

\clearpage
\section*{5. Пример: ограничение на каждый платёж}

Кредит $S$ (в млн руб.) выдан под $20\%$ годовых на три года. Остатки после первого, второго и третьего платежей: $0{,}8S$, $0{,}4S$, $0$.

\[
\begin{aligned}
P_1&=1{,}2S-0{,}8S=0{,}4S,\\
P_2&=1{,}2\cdot0{,}8S-0{,}4S=0{,}56S,\\
P_3&=1{,}2\cdot0{,}4S=0{,}48S.
\end{aligned}
\]

Требование «каждый платёж меньше $5$ млн руб.» записывается системой:
\[
\begin{cases}
0{,}4S<5,\\
0{,}56S<5,\\
0{,}48S<5.
\end{cases}
\quad\Longleftrightarrow\quad
0<S<\frac{125}{14}\approx8{,}93.
\]

\key{Внимание к строгости знака!} Если $S$ --- произвольное положительное действительное число, наибольшего значения нет: верхняя граница не достигается. Если по условию $S$ обязано быть \emph{целым числом миллионов рублей}, наибольшее значение $S=8$.

\section*{6. Смешанная схема: сначала только проценты}

Кредит $S=3{,}3\mln$ взят на пять лет под $20\%$ годовых. После платежей первых трёх лет долг остаётся равным $S$. В четвёртом и пятом годах платежи равны; к концу пятого года кредит погашен.

Пусть $k=1{,}2$, $r=0{,}2$. В первые три года выплачиваются \key{только проценты}: $P_1=P_2=P_3=rS=0{,}66\mln$. Поэтому основной долг остаётся равным $S$.

Два последних платежа обозначим $X$. Для четвёртого и пятого годов:
\[
\begin{aligned}
B_4&=kS-X,\\
B_5&=k(kS-X)-X=k^2S-(k+1)X=0.
\end{aligned}
\]

Следовательно,
\[
X=\frac{k^2S}{k+1}
=\frac{1{,}2^2\cdot3{,}3}{2{,}2}=2{,}16\mln.
\]

Разница последнего и первого платежей:
\fboxeq{2{,}16-0{,}66=1{,}50\mln.}

\textbf{Вывод:} схема смешанная; равны только два последних платежа.

\clearpage
\section*{Алгоритм: произвольная кредитная схема}
\vspace{8mm}
\begin{center}
\begin{tikzpicture}[
  >=Latex,
  start/.style={ellipse,draw=accent,very thick,minimum width=38mm,minimum height=10mm,align=center,fill=accent!5},
  action/.style={rectangle,rounded corners=2pt,draw=accent,thick,text width=108mm,minimum height=11mm,align=center,inner sep=5pt},
  branch/.style={diamond,aspect=2.9,draw=darkaccent,thick,align=center,text width=39mm,inner sep=2pt},
  split/.style={rectangle,rounded corners=2pt,draw=accent,thick,text width=60mm,minimum height=18mm,align=center,inner sep=4pt},
  arr/.style={-{Latex[length=2.4mm]},thick,color=darkaccent},
  font=\small
]
\node[start] (begin) at (0,0) {Начало};
\node[action] (data) at (0,-1.6) {Определите $S$, ставку за период $r$ и даты начислений и платежей};
\node[action] (balances) at (0,-3.2) {Выпишите $B_0=S$ и все известные остатки долга};
\node[action] (interest) at (0,-4.8) {Для очередного периода: $I_i=rB_{i-1}$, $D_i=(1+r)B_{i-1}$};
\node[branch] (known) at (0,-7.0) {Известен остаток $B_i$?};
\node[split] (yes) at (-4.05,-9.5) {Найдите платёж:\par $P_i=D_i-B_i$};
\node[split] (no) at (4.05,-9.5) {Обозначьте платёж $P_i$ (при равенстве --- $X$);\par $B_i=D_i-P_i$};
\node[action] (next) at (0,-12.0) {Повторите расчёт для остальных периодов; используйте $B_n=0$, если кредит погашен};
\node[action] (condition) at (0,-13.9) {Составьте требуемое уравнение или неравенство; проверьте единицы измерения и строгие границы};
\node[start] (end) at (0,-15.6) {Ответ и проверка};
\draw[arr] (begin)--(data);
\draw[arr] (data)--(balances);
\draw[arr] (balances)--(interest);
\draw[arr] (interest)--(known);
\draw[arr] (known.west) -| node[pos=0.20,above] {да} (yes.north);
\draw[arr] (known.east) -| node[pos=0.20,above] {нет} (no.north);
\coordinate (join) at (0,-11.1);
\draw[thick,color=darkaccent] (yes.south) |- (join);
\draw[thick,color=darkaccent] (no.south) |- (join);
\draw[arr] (join)--(next.north);
\draw[arr] (next)--(condition);
\draw[arr] (condition)--(end);
\end{tikzpicture}
\end{center}

\clearpage
\section*{7. Тренировочный блок --- Путь А}
\textbf{10 задач по нарастающей сложности.} В каждом задании проценты начисляются \emph{перед} очередным платежом; если указана ставка $p\%$, она действует за один соответствующий период.

\begin{enumerate}
\item В начале месяца долг составлял $200\,000\rub$. За месяц он вырос на $5\%$. После платежа осталось $180\,000\rub$. Найдите платёж.

\item В начале года долг составлял $400\,000\rub$. Банк начислил $10\%$, затем заёмщик заплатил $170\,000\rub$. Каков остаток долга?

\item Кредит $1\mln$ выдан на три года под $20\%$ годовых. После платежей первого, второго и третьего годов остатки составляют соответственно $0{,}7S$, $0{,}4S$ и $0$. Найдите все три платежа в миллионах рублей.

\item Кредит $S$ выдан на три месяца под $4\%$ в месяц. Остатки после первого, второго и третьего платежей равны $0{,}75S$, $0{,}25S$, $0$. На сколько процентов общая сумма платежей больше суммы кредита?

\item В начале года долг составлял $800\,000\rub$. За год банк начислил $15\%$. После платежа долг стал на $30\%$ меньше суммы, имевшейся \emph{в начале года}. Найдите платёж.

\item Кредит $S$ миллионов рублей выдан на три года под $20\%$ годовых. Остатки после платежей равны $0{,}8S$, $0{,}4S$, $0$. Каждый платёж строго меньше $5\mln$. Найдите наибольшее $S$, если $S$ --- положительное целое число миллионов рублей.

\item Кредит $S$ миллионов рублей выдан на три года под $10\%$ годовых. Остатки после платежей равны $0{,}6S$, $0{,}3S$, $0$. Общая сумма платежей строго меньше $15\mln$. Найдите наибольшее $S$, если $S$ --- положительное целое число миллионов рублей.

\item Кредит $2{,}2\mln$ выдан на четыре года под $20\%$ годовых. В первые два года остаток долга после каждого платежа равен первоначальной сумме кредита. В третьем и четвёртом годах платежи одинаковы; к концу четвёртого года долг равен нулю. На сколько миллионов рублей последний платёж больше первого?

\item Кредит $S$ выдан на два года под $25\%$ годовых. После первого платежа осталось $0{,}6S$, после второго --- $0$. Общая сумма двух платежей равна $7\mln$. Найдите первоначальную сумму кредита.

\item Кредит $S$ выдан на три года под $20\%$ годовых. В конце первого года остаток долга стал на $70\%$ меньше начального, в конце второго --- на $70\%$ меньше остатка после первого года, а в конце третьего долг выплачен полностью. На сколько процентов общая сумма платежей больше первоначального кредита?
\end{enumerate}

\clearpage
\section*{Ответы к тренировочному блоку}
\begin{center}
\renewcommand{\arraystretch}{1.65}
\begin{tabularx}{\linewidth}{@{}c >{\raggedright\arraybackslash}X@{}}
\toprule
№ & Ответ \\
\midrule
1 & $30\,000\rub$ \\
2 & $270\,000\rub$ \\
3 & $P_1=0{,}50\mln$; $P_2=0{,}44\mln$; $P_3=0{,}48\mln$ \\
4 & $8\%$ \\
5 & $360\,000\rub$ \\
6 & $8\mln$ \\
7 & $12\mln$ \\
8 & $1\mln$ \\
9 & $5\mln$ \\
10 & $27{,}8\%$ \\
\bottomrule
\end{tabularx}
\end{center}

\vspace{7mm}
\key{Контрольные вопросы для самостоятельной проверки}
\begin{itemize}
\item Учитывали ли Вы проценты до каждого платежа?
\item Отличали ли Вы условие на \emph{каждый} платёж от условия на \emph{сумму} платежей?
\item Проверили ли Вы смысл выражений «на $70\%$ меньше» и «равен $70\%$»?
\item Если использовалось строгое неравенство, проверили ли Вы, достигается ли найденная граница?
\end{itemize}

\end{document}
